\scanpage{141}
\mspage{70}

The image of this map thus consists of the \inserted{$\pi$-conjugacy classes of the} section subgroups $\Gamma'\subset\mathcal E$ such that $\widetilde U^{\Gamma'}\ne\varnothing$, and for such a $\Gamma'$, the set of \unsure{$x\in U^\Gamma$} which \deleted{\unsure{to it}} give as image this conjugacy class is the image of $U^{\Gamma'}$ in $U$. Finally, if $x$ is in this image, the set \inserted{$\widetilde U^{\Gamma'}_x$} of the $\widetilde x\in\widetilde U^{\Gamma'}$ over $x$ (\deleted{\unsure{$\ne\varnothing$}} hypothesis on \unsure{a|x}) is a torsor under $\pi^{\Gamma'}$ i.e. if $\widetilde x\in(\widetilde U^{\Gamma'})_x$, \deleted{\unsure{of}} and if $\alpha\in\pi$, then
\[
\alpha.\widetilde x\in\widetilde U^{\Gamma'}\quad\Longleftrightarrow\quad\alpha\in\pi^{\Gamma'}
\]
(Trivial verification, as for that matter all the preceding assertions). Thus, the following conditions are equivalent:

\margin{N.B. \unsure{$\widetilde U_x$} being identified with \unsure{$\operatorname{Isom}_U(R_x,\widetilde U)$}, $\widetilde U_x^{\Gamma'}$ is identified with \unsure{$\operatorname{Isom}_{U,\Gamma}(R_x,\widetilde U_{\Gamma'})$} \unsure{of the two isomorphisms that commute|of covering isomorphisms commuting} with the action of $\Gamma$ and $\pi^\Gamma=(1)$ means that this isomorphism is unique \unsure{i.e.} \deleted{\illegible{3}} \unsure{the conjugacy class of the sections} \illegible{2}.}

 a) $(*)$ is a bijection, and $\forall$ subgroup $\Gamma'$ of $\mathcal E$ a section of the extension, \deleted{\illegible{1}} we have $\pi^{\Gamma'}=\{1\}$.

 b) $\forall\,\Gamma'$ as in a), $\Gamma'$ operating on $\widetilde U$ has a fixed poin\textsuperscript{t} and only one.

This being so, let us prove the

\uline{Fundamental lemma} Let $\Gamma$ be a \inserted{finite} finite group, operating faithfully on \inserted{a space} \unsure{$D$}\,$\simeq\mathbb R^2$, let \deleted{We assume $\Gamma=$} $\Gamma^\circ$ be the subgroup of $\Gamma$

\mspage{140}

\scanpage{142}

(of index $0$ or $1$) consisting of the $g\in\Gamma$ such that $g_D$ preserves the orientation, and suppose $\Gamma^\circ\ne\{1\}$ (i.e. $\Gamma$ is reduced neither to the trivial group, nor to the group \deleted{\illegible{1}} $\{1,\sigma\}$, where $\sigma$ is an anti-involution of $D$. Then

 a) $\Gamma$ has a fixed poin\textsuperscript{t} and only one in $D$ i.e. $\operatorname{card}D^\Gamma=1$

 b) \deleted{\unsure{And}} $\Gamma^\circ$ is cyclic, and if $\Gamma\ne\Gamma^\circ$, $\Gamma$ is a \inserted{dihedral} group.

More precisely, there exists a homeomorphism
\[
D\simeq\mathbb C
\]
such that \deleted{$\Gamma$ operates by automorphisms} \inserted{\unsure{on $\mathbb C$ by the group}} \deleted{\illegible{1}} of homeomorphisms of $\mathbb C$ transformed from \deleted{complex and anti-complex of $\mathbb C$} \deleted{\unsure{fixing}} $\Gamma$ is either the group of homotheties by $\mu_n(\mathbb C)$ (if $\Gamma=\Gamma^\circ$ of order $n$), or the associated dihedral group \unsure{$z\longrightarrow\zeta\tau^\varepsilon(z)$} ($\zeta\in\mu_n(\mathbb C)$, $\varepsilon=\pm1$) where \unsure{$\sigma$} \deleted{\illegible{1}} is complex conjugation.

\uline{Proof of the fundamental lemma.}

\deleted{Suppose first $\Gamma=\Gamma^\circ$, i.e. $\Gamma$ preserves the orientation. \deleted{If $\Gamma$ operated \unsure{without a poin\textsuperscript{t}}} We proceed}

\mspage{141}

\scanpage{143}
\mspage{71}

\deleted{By induction on the order of $\Gamma$.}

\deleted{\unsure{a)} \unsure{$D^\Gamma\ne\varnothing$}}

\deleted{a) Case $\Gamma$ of prime order}

 a) Suppose first that one can find a $C^2$ structure on $D$ invariant under $\Gamma$ --- then a standard argument shows that there exists a conformal structure invariant under $\Gamma$, now the fundamental theorem of conformal mapping shows that then
\[
\begin{aligned}
&\text{either }D\simeq\text{interior}\ \inserted{\Delta}\ \text{of the unit disk }\Delta\ \inserted{\text{or of the Poincaré half-plane}}\\
&\text{or }D\simeq\mathbb C\qquad\text{(conformal isomorphisms).}
\end{aligned}
\]
In the first case, the group of conformal automorphisms \deleted{maps} is
\[
\simeq\unsure{\operatorname{Sl}(2,\mathbb R)^\sim\,|\,\operatorname{Gl}(2,\mathbb R)^\sim}
=\{u\in\operatorname{Gl}(2,\mathbb R)\mid\det u=\pm1\}
\]
(\deleted{operating by} $u=\left(\begin{smallmatrix}a&b\\c&d\end{smallmatrix}\right)$ operating on \deleted{\illegible{1} Poincaré} \inserted{$\widehat{\mathbb C}$} by \deleted{\unsure{$z\longmapsto\dfrac{az+b}{cz+d}$}} $\Theta_u\circ\tau^{\det u}$, where $\tau$ is complex conjugation and $\Theta_u(z)=\dfrac{az+b}{cz+d}$, and leaving the Poincaré half-plane stable), and every compact subgroup is \inserted{contained in a} conjugate of the maximal compact \uline{subgroup}, which (on passing back to the unit disk \inserted{$\Delta$}) is identified with the group $O(2,\mathbb R)$ of

\mspage{142}

\scanpage{144}

transformations of the unit disk of the form
\[
z\longrightarrow\zeta\tau^\varepsilon(z)\qquad
\begin{gathered}
\zeta\in\mathbb U=\{\zeta\in\mathbb C\mid |\zeta|=1\}\\
\varepsilon\in\{\pm1\}\\
\tau\ \text{complex conjugation}
\end{gathered}
\]
Every finite subgroup of this group is of one of the types made explicit above --- we win, by using a \unsure{homeomorphism} $[0,1[\longrightarrow[0,+\infty[$ to define a homeomorphism \deleted{\illegible{1}} $\Delta\simeq\mathbb C$ commuting with the action of $G=\deleted{\unsure{S}}O(2,\mathbb R)$.

In the case $D\simeq\mathbb C$, we see that \deleted{\unsure{is}} the group of conformal automorphisms of $\mathbb C$ is the group of transformations $az+b$ or $a\overline z+b$, in which a maximal compact subgroup is the same group $\deleted{\unsure{S}}O(2,\mathbb R)$ as just now --- and every compact subgroup (a fortiori, every finite subgroup) is contained in a conjugate of this one. We win again.

The rest of the work consists \inserted{essential\textsuperscript{ly}} in \deleted{eliminating} showing that the hypothesis of non-wildness

\mspage{143}

\scanpage{145}
\mspage{72}

is always satisfied, \inserted{at least for $\Gamma^\circ$. Suppose first $\Gamma=\Gamma^\circ$.} \deleted{We assume} \inserted{We assume} that all is proved for \deleted{the groups} \inserted{the orders} $<\operatorname{card}(\Gamma)$. \deleted{Suppose}

 b) $\Gamma$ has a fixed poin\textsuperscript{t} \inserted{\unsure{$o$|$0$}}, i.e. $D^\Gamma\ne\varnothing$.

Otherwise, the stabilizer subgroups \inserted{$\Gamma_x$} of the $x\in D$ being of order $<$ \deleted{\illegible{1}} $\operatorname{card}\Gamma$, by the induction hypothesis, the $\Gamma_x$ have the structure stated in the theorem, hence \unsure{$D^\Gamma=U$|$D/\Gamma=U$} is a topological surface and \unsure{$D\longrightarrow U$|$D\longrightarrow D^\Gamma$} is a ramified covering, let us choose a \deleted{complex conformal} \inserted{conformal} structure on $U$ \deleted{N.B. $U$ is oriented}; there is a unique conformal structure on $D$ such that $D\longrightarrow U$ is \deleted{\unsure{conformal|holomorphic}} \inserted{``conformal''} (holomorphic or anti-holomorphic), this is invariant under $\Gamma$, and by a) $\Gamma$ has a fixed poin\textsuperscript{t}, contradiction.

 c) $\Gamma$ has no fixed poin\textsuperscript{t} other than \unsure{$o$|$0$}. We let $\Gamma$ operate \inserted{\unsure{faithful\textsuperscript{ly}}} on \unsure{$D^*=D\setminus\{o\}$|$D^*=D\setminus\{0\}$}, and we have to prove that $D^{*\Gamma}=\varnothing$.

\mspage{144}

\scanpage{146}

Let then $x\in D^{*\Gamma}$ --- we are going to reach a contradiction. Consider the universal covering $\widetilde D^*$ of $D^*$ pointed at $x$, \deleted{We may assume} so $\Gamma$ operates on $\widetilde D^*$ \unsure{with} a fixed poin\textsuperscript{t} $\widetilde x$ over $x$. Here $\pi=\pi_1(D^*)\simeq\mathbb Z$, and $\Gamma$ operates on it \unsure{trivial\textsuperscript{ly}} (since $\Gamma=\Gamma^\circ$) so $\Gamma\times\mathbb Z$ operates on $\widetilde D^*$. We may assume $D=\mathbb C$, \unsure{$o=0$|$0=0$}, $x=1$, $D^*=\mathbb C^*$, $\widetilde D^*=\mathbb C$, $\widetilde x=0$, $\widetilde D^*\longrightarrow\widetilde D$ given by $\exp 2i\pi z$, \deleted{\unsure{so}} and $\mathbb Z$ operating on $\mathbb C$ by $\Theta_nz=z+n$ ($n\in\mathbb Z$).

\deleted{Lemma} It remains to prove that if a \inserted{finite group} $\Gamma$ operates on $\mathbb C$ commuting \deleted{\unsure{with}} with the action of $\mathbb Z$ on $\mathbb C$, and leaving \deleted{a poin\textsuperscript{t}} fixed the poin\textsuperscript{t} \unsure{$o$|$0$}, then $\Gamma$ operates trivially (which contradicts the hypothesis of faithfulness of the operation).

\deleted{So} We are reduced to the.

\uline{Lemma} Let $u$ be a homeomorphism \inserted{of finite order} of $\mathbb C$,

\mspage{145}

\scanpage{147}
\mspage{73}

commuting with $z\longmapsto z+1$, and leaving the origin invariant. Then $u=\operatorname{id}$

We may assume \deleted{$u$ of prime order} $\exists$ prime number $p$ such that $u^p=\operatorname{id}$, i.e. that $u$ corresponds to an operation of $\Gamma=\mathbb Z/p\mathbb Z$ on $\mathbb C$. \inserted{All the poin\textsuperscript{ts} of $\mathbb Z\subset\mathbb C$ are fixed by $\Gamma$.} Passing to $\widehat{\mathbb C}\simeq S^2$, we find that $\infty$ is an accumulation poin\textsuperscript{t} of the fixed poin\textsuperscript{ts} under $\Gamma$ --- if $\Gamma$ did not operate trivially, this would be decidedly very wild! We ought to be able to finish by the Adams spectral sequence... I am not going into these mazes...

 d) The purely topological part being \unsure{admitted|thus} \unsure{assumed|assumed} \unsure{proved|proved}, we thus conclude from this, if $\Gamma\ne\Gamma^\circ$, $\Gamma^\circ\ne(1)$, \deleted{that the unique} \inserted{as} \deleted{fixed poin\textsuperscript{t} of} \inserted{$D^{\Gamma^\circ}$} is invariant under $\Gamma$, \deleted{that} \inserted{as} $D^{\Gamma^\circ}$ is reduced to a poin\textsuperscript{t}, this one is invariant under all of $\Gamma$, not just $\Gamma^\circ$. On the other hand, We know enough now to know that if $\Gamma$ finite group operates on a surface \deleted{Hausdorff} \inserted{paracompact $U$}, the $\Gamma_x$ ($x\in U$) respecting the orientation, then \unsure{$U/\Gamma=V$}

\mspage{146}

\scanpage{148}

is a surface --- \inserted{and $U\longrightarrow V=U/\Gamma$ is a ramified covering} and choosing a conformal structure on this one, we find by inverse image a conformal structure on $U$ invariant under $\Gamma$, \deleted{and} \inserted{For the case $U=D$} we finish by a), \deleted{but} for the complement to the fundamental lemma. [But to do things properly, one would have to prove that there always exists invariant conformal structure, \deleted{\illegible{2}} $\Gamma$ finite group operates on surfaces \deleted{\illegible{1}} \unsure{paracompact} --- so $U/\Gamma$ is a surface with boundary.... Here what is missing is the analysis of the action of an anti-involution of a surface in a neighborhood of a fixed poin\textsuperscript{t}...]

\margin{One would have to prove that if $\tau$ is an involutive \unsure{anti-automorphism} of $D$, then $\exists$ isomorphism $D\simeq\mathbb C$ such that $\tau$ becomes $z\longmapsto\overline z$ (hence $D^\tau\simeq\mathbb R$!) which should make it possible to prove, if $\Gamma\simeq\mathbb Z/2\mathbb Z$ operates by anti-automorphisms on $U$ (\unsure{oriented|orientable} $U\ne S^2$) that $\pi_0(U^\Gamma)$ \unsure{$\Longrightarrow$} \unsure{$\pi$-conjugacy classes of sections of $\mathcal E$ over $\Gamma$} is bijective.}

Consequence of the fundamental lemma:

\uline{Theorem.} Let $U$ be a paracompact topological surface \inserted{$0$-connected}, $\Gamma$ a finite group operating \uline{faithfully} on $U$, we assume $\widetilde U$ \deleted{\illegible{1}} noncompact (i.e. $U$ not homeomorphic to $S^2$ nor

\mspage{147}

\scanpage{149}
\mspage{74}

to the real projective plane), we assume that \deleted{\unsure{$\Gamma\ne(1)$}, \illegible{1}} that \inserted{moreover} if $U$ is \deleted{not} orientable, then \deleted{\unsure{$\Gamma\ne\{\pm1\}$} i.e. $\operatorname{card}\Gamma\ge3$ \illegible{1}} \deleted{we} assume that the subgroup $\Gamma^\circ$ of $\Gamma$ of the $g\in\Gamma$ which preserve an orientation is $\ne(1)$ [hence $\Gamma$ \unsure{not reduced} either to $1$, or to \unsure{$1$ and|Id} an anti-involution], and if $U$ nonorientable, that \deleted{\unsure{cardinality}} $\Gamma$ \unsure{$\not\simeq$|$\ne$} $(1),\{\pm1\}$ i.e. $\operatorname{card}\Gamma\ge3$. This being so, \deleted{\unsure{the map}} let us consider \deleted{the map} \inserted{the extension} $\mathcal E$ of $\Gamma$ by $\pi=\pi_1(U)$, and the map
\[
U^\Gamma\longrightarrow\begin{gathered}
\text{Classes of }\pi\text{-conjugacy of}\\
\text{splittings of }\mathcal E\longrightarrow\Gamma
\end{gathered}
\]
we have this:

 a) This map is bijective

 b) For every section subgroup $\Gamma'\subset\mathcal E$, we have $\pi^{\Gamma'}=(1)$

 c) If $U^\Gamma\ne\varnothing$ i.e. If there exists a splitting, then $\Gamma$ is cyclic or dihedral (and in the case $U$ oriented, $\Gamma^\circ$ is cyclic)

\uline{N.B.} The hypotheses on $U$ ensure that $\widetilde U\simeq D$, and those on $\Gamma$ that $\Gamma$ operating on $D$ satisfies the hypotheses of the fundamental lemma.

\mspage{148}

\scanpage{150}

\deleted{\uline{Corollary-Scholium}, For an $x\in U^\Gamma$ corresponding to a splitting class of the extension $\mathcal E$, we find that the splittings of this class, \deleted{on which $\mathcal E$ operates by \unsure{inner automorphisms|functoriality}}, form a \deleted{\unsure{$\mathcal E$-set which is a}} $\pi$-torsor, \deleted{and the automorphisms of}}

\uline{Corollary} Suppose \deleted{that} $U$ \unsure{orientable}, \deleted{\illegible{1}} $\Gamma$ preserving the orientation. Let $U^!$ be the set of the $x\in U$ such that $\Gamma_x\ne(1)$. (which is therefore a discrete subset of $U$).

To every $x\in U^!$, let us associate the conjugacy class of the subgroups of $\mathcal E$ \inserted{partial \inserted{\unsure{$\pi$}} sections of $\mathcal E$ over $\Gamma_x$} which \unsure{correspond|corresponding} to \unsure{this|this} $x\in U^{\Gamma_x}$. Then

 a) \deleted{the map} the partial sections thus obtained are \uline{maximal} among those which are $\ne(1)$

 b) The map from $U^!$ to \unsure{the set|one} of the $\pi$-conjugacy classes of maximal partial sections \deleted{nontrivial} \inserted{$\ne(1)$}, is bijective

\margin{\uline{N.B.} This map commutes with the natural actions of $\Gamma$!}

 c) For every such partial section, we have $\pi^{\Gamma'}=(1)$, i.e. the \deleted{\unsure{$\Gamma_x$-}} \unsure{automorphisms} of the \inserted{universal} covering \unsure{$R_x$} of $U$ which commute with the action of $\Gamma_x$ are trivial

--- there is a unique isomorphism \inserted{commuting with the action of $\Gamma_x$} between this torsor, and the torsor deduced from $\widetilde U$ by twisting by the $\pi$-torsor \deleted{\unsure{of the}} $P_x$ of the \unsure{$\Gamma''$} in the class \deleted{\illegible{1}} \unsure{$[\Gamma']$}...

\mspage{149}